\chapter{Status of the Open Problems}
\label{app:status}

This appendix records the precise status of Problems A--J as of the present text.
It collects standard supporting calculations that are already rigorous, states
the exact missing step for each problem, and does not convert any open problem
into a lemma. The research agenda of Chapter~\ref{ch:agenda} remains the
authoritative list; this appendix is a working ledger.

\section{Problem A: matrix-to-modular map}
\label{sec:status-A}

The Weyl--Wigner expansion
\begin{equation}
  [\hat{X}_\mu,\hat{\Phi}]
  =i\Theta_{\mu\nu}D^\nu\Phi
  +\frac{1}{24}\Theta_{\mu\alpha}\Theta_{\beta\gamma}D^\alpha D^\beta D^\gamma\Phi
  +\mathcal{O}(\Theta^3)
\end{equation}
is a standard result of the Moyal product (see the infrared dictionary
discussion in Section~\ref{sec:ch3-relate}). This expansion relates the matrix commutator to a continuum
derivative, but it does not provide the required map from
\((\hat X,\Phi)\) to the causal-diamond modular Hamiltonian \(H_D\)
or the mutual-information functional \(U_{\mathrm{ent}}[I(A:B)]\)
before the composite metric is defined.

\paragraph{Missing step.}
An error-controlled map from the matrix order parameter (or from
\([\hat{X},\Phi]\)) into the modular Hamiltonian \(H_D\) of a causal diamond,
or into the mutual-information functional \(U_{\mathrm{ent}}[I(A:B)]\),
that holds before the composite metric is defined. Until that map is a lemma,
Chapter~8 computes a heat kernel in a theory whose relation to Chapter~3 is
still Assumption~\ref{ass:dictionary}.

\section{Problem B: replica Ward identity}
\label{sec:status-B}

For a differentiable family of normalized reduced states on a fixed
diamond algebra, the first variation of the entropy obeys the standard
first-law identity
\[
\delta S_D
=
\operatorname{Tr}(\delta\rho_D H_D^{(0)})
=
\delta\langle H_D^{(0)}\rangle,
\qquad
H_D^{(0)}=-\ln\rho_D^{(0)},
\]
where the modular Hamiltonian is held fixed at the reference state.
Under the usual path-integral assumptions, the replica construction
expresses this variation as a response of the microscopic action
insertion on the branched cover. These identities do not imply
\(\delta S_D=0\).

\paragraph{Missing step.}
Derive from the microscopic dynamics the entanglement-stationarity
condition \(\delta S_D=0\) for the allowed fixed-volume variations,
and establish the area contribution and its relation to the
\emph{emergent} metric without presupposing that metric or the
fluctuation split of Chapter~5. For variations of the auxiliary
scaffold, the replica insertion computes a response to
\(g^{(0)}_{\mu\nu}\); relating it to the emergent metric still
requires a specified map between them. Problem~B therefore remains
open, although the fixed-algebra first-law identity is established.

\section{Problem C: conformal mode and Cassini}
\label{sec:status-C}

The radial singlet has the tree-level mass
\[
m_\sigma^2=\lambda v^2.
\]
This does not by itself determine the mass or matter coupling of a
scalar mode of the composite metric. In particular, the stated
Repair~I prescription is intended to build a vierbein from the Lorentz-valued
polar factor \(\Lambda\), evaluated at \(\sigma=0\), rather than from
the radial field.

\paragraph{Obstruction to the stated Repair I frame.}
Suppose \(\Lambda(x)\in\mathrm{SO}(3,1)\) and expand about the locked
configuration,
\[
\Lambda=\mathbb I+\Omega+O(\Omega^2),
\qquad
\Omega_{\mu\nu}=-\Omega_{\nu\mu}.
\]
Then \(\partial_\mu\Omega_{\alpha\beta}\) transforms in
\(V\otimes\mathfrak{so}(3,1)\), where \(V=(\tfrac12,\tfrac12)\)
is the Lorentz vector representation. After complexification,
\[
\bigl(V\otimes\mathfrak{so}(3,1)\bigr)_{\mathbb C}
=
\left(\tfrac32,\tfrac12\right)
\oplus\left(\tfrac12,\tfrac12\right)
\oplus\left(\tfrac12,\tfrac32\right)
\oplus\left(\tfrac12,\tfrac12\right).
\]
A vierbein fluctuation transforms in \(V\otimes V\), whose
complexified decomposition is
\[
(V\otimes V)_{\mathbb C}
=
(0,0)\oplus(1,0)\oplus(0,1)\oplus(1,1).
\]
These representations have no common irreducible summand. Hence
there is no nonzero Lorentz-equivariant linear map from
\(\partial_\mu\Omega_{\alpha\beta}\) to a vierbein fluctuation using
only the locked Lorentz-invariant tensors. Consistently, the displayed
formula
\[
e^A{}_\mu=\left\langle\partial_\mu\Lambda^A{}_I\right\rangle
\]
leaves \(I\) uncontracted and does not, as written, define a vierbein.

A simple proposed contraction with a fixed nonzero internal vector
\(n^I\) does not repair non-degeneracy. If
\[
e^A{}_\mu=\partial_\mu\Lambda^A{}_I\,n^I,
\]
then Lorentz orthogonality gives
\[
(\Lambda n)_A e^A{}_\mu
=
(\Lambda n)_A\partial_\mu(\Lambda n)^A
=
\frac12\partial_\mu\bigl((\Lambda n)^2\bigr)
=0.
\]
Every column of \(e^A{}_\mu\) therefore lies in the
three-dimensional subspace orthogonal to \(\Lambda n\), so
\(\operatorname{rank}(e)\leq 3\). This contraction cannot produce a
non-degenerate four-dimensional frame.

\paragraph{Status.}
Under the stated Repair~I ingredients, the linearized composite
vierbein required for the later metric-scalar analysis is not
constructed. This is an obstruction to that frame prescription, not
a proof that a completed theory has or lacks a physical scalar pole.
Repair~II, an independent coframe, or additional symmetry-breaking
background tensors would change the construction and require a
separate analysis. The radial mass \(m_\sigma^2=\lambda v^2\) remains
established, but it is not a metric-scalar mass. The scalar spectrum
and matter coupling relevant to Cassini are therefore still
undetermined. The Repair~II pure-trace relation \(\psi=8\sigma/v\)
with \(m_\sigma^2=\lambda v^2\) is a potential-only statement; the
induced Einstein--Hilbert term adds to the kinetic coefficient of
\(\sigma\), so it does not fix a physical pole mass, mixing or
Cassini coupling.

\section{Problem D: local Lorentz}
\label{sec:status-D}

The counting of generators under the locked vacuum is given in
Chapter~10. The condensate term supplies a quadratic Higgs-type
contribution in the broken-direction combination \(\Xi\); it does not
establish a physical pole mass without kinetic terms and canonical
normalization. The diagonal connection's dynamics and torsion remain
open. Torsion-free behavior follows only under the conditional
Palatini--Holst hypotheses of Proposition~
\ref{prop:B-torsion-free}; NAMCG has not been shown to induce those
hypotheses.

For a Laplace-type operator,
\(a_1=\operatorname{tr}(R(g)/6-E)\); the explicit universal
curvature term is the Levi-Civita scalar \(R(g)\). Because \(E\) is
unspecified, no conclusion about additional \(B\)-dependence of
\(a_1\) follows, and a Palatini \(R(B)\) term is not established. If
\(B\) enters the bundle connection of the actual operator, \(a_2\)
may contain a \(\operatorname{tr}F(B)^2\) term; whether it survives
in the weighted trace depends on the unspecified operator
representation and weights. Deriving the \(B\)-dependence of \(E\)
and the bundle connection remains open.

\paragraph{Missing step.}
Derive \(\mathrm{SO}(3,1)\) from the matrix algebra alone, or state
that Presentation~A is not ontologically prior (the present text
does the latter), and derive the effective dynamics of the
independent diagonal connection or retain it as unresolved field
content. Generator counting alone does not settle its torsion or
propagating modes.

\section{Problem E: signature}
\label{sec:status-E}

Euclidean positivity of the matrix action and the formal Wick rotation of
one coordinate are recorded. They do not select Lorentzian signature without
an additional hypothesis.

\paragraph{Status.}
Remains Assumption~\ref{ass:lorentzian}.

\section{Problem F: coefficients \(A,B,C\)}
\label{sec:status-F}

The conditional, truncated heat-kernel identity is
Theorem~\ref{thm:Geff}; it is not a computed \(A,B,C\)
parameterization or physical value of \(G_{\mathrm{eff}}\).

\paragraph{Missing step.}
Specify the gauge-fixed fluctuation, Goldstone, diffeomorphism-ghost,
and local-Lorentz-ghost operators and weights; per-sector proper-time
limits; massive threshold functions if masses are resummed; the
continuation; and the finite-counterterm scheme. Then evaluate the
weighted \(a_1\) and curvature-linear \(a_2\) contributions in that
scheme. Until this is completed, Problem~F remains open, and the
asserted \(N_{\mathrm{eff}}=+2\) count does not close it.

\section{Problem G: HBEV}
\label{sec:status-G}

The boundary functional of Chapter~12 regulates \(\Lambda_{\mathrm{eff}}\)
once an infrared entropy density is supplied. The working value
\(\Lambda_{\mathrm{eff}}=3/L_{\mathrm{IR}}^2\) is a choice consistent with
that functional, not a derived necessity.

\paragraph{Status.}
Closed as a consistent regulation mechanism under the stated boundary
entropy; not closed as a derivation of the entropy density from the matrix
theory.

\section{Problem H: modular origin of the interpolating function}
\label{sec:status-H}

Appendix~\ref{app:empirical} tests the function
\[
\nu_{\sqrt{}}(g_{\mathrm N};g_s)
=
\frac{1}{1-\exp\!\left(-\sqrt{g_{\mathrm N}/g_s}\right)}
\]
against SPARC with \(g_s\) fixed to the empirical value
\(g_\dagger\). This is a coefficient-level consistency check, not
a microscopic derivation. Appendix~\ref{app:matrix_entropy} gives
the standard heat-kernel area contribution but does not compute a
source-dependent modular spectrum. Appendix~\ref{app:gaia} applies
an external-field prescription to wide binaries; it does not derive
that prescription from the modular data.

\paragraph{Conditional lemma: occupation of one modular ladder.}
Suppose the relevant reduced state contains one bosonic ladder with
normalized probabilities
\[
p_n=(1-e^{-\Delta k})e^{-n\Delta k},
\qquad n=0,1,2,\ldots,\qquad \Delta k>0.
\]
Then its mean occupation is
\[
\langle n\rangle
=
\sum_{n=0}^{\infty} n p_n
=
\frac{1}{e^{\Delta k}-1},
\]
and therefore
\[
1+\langle n\rangle
=
\frac{1}{1-e^{-\Delta k}}.
\]
Thus, if the extra-stress response is exactly \(1+\langle n\rangle\)
for this single mode, the target interpolator follows if and only if
\[
\Delta k(g_{\mathrm N})
=
\sqrt{\frac{g_{\mathrm N}}{g_s}}.
\]
The quantity to calculate is the modular excitation gap
\(\Delta k\), not the absolute lowest eigenvalue of \(H_D\), which
depends on an additive normalization constant.

For a multimode response, the corresponding quantity is instead a
weighted spectral sum, schematically
\[
\mathcal R(g_{\mathrm N})
=
\int_{(0,\infty)}
\frac{d\mu_{g_{\mathrm N}}(\Delta k)}
     {e^{\Delta k}-1},
\]
where \(d\mu_{g_{\mathrm N}}\) includes the mode weights coupling
to the extra stress. In that case a square-root gap for one mode
does not by itself establish the target function; the measure and
its source dependence must also be calculated.

\paragraph{Missing step.}
Construct the regulated modular operator, or an equivalent
well-defined spectral response, for a weak-field causal diamond
containing a baryonic source. Derive the relevant excitation gap
\(\Delta k(g_{\mathrm N},\kappa_{\mathrm{IR}})\) and its spectral
weight from the microscopic theory. Fix the normalization relating
\(g_s\) to \(\kappa_{\mathrm{IR}}\) independently of galaxy data.
In particular, reconcile the convention
\(g_s\sim\kappa_{\mathrm{IR}}/(2\pi)\) in
Appendix~\ref{app:rar-bridge} with the undetermined factor
\(g_s=c_\star\kappa_{\mathrm{IR}}\) in Chapter~17 before claiming
that the Bose occupation yields the target.

Until this calculation is supplied, Appendices~\ref{app:empirical},
\ref{app:matrix_entropy}, and~\ref{app:gaia} provide an empirical fit
check, a standard area-term calculation, and an external-field
consistency test, respectively; together they do not derive or
uniquely select \(\nu_{\sqrt{}}\). Problem~H remains open.

\section{Problem I: dynamical origin of the drag}
\label{sec:status-I}

Chapter~\ref{ch:dissipation} adopts a vector drag proportional to
\(-\mathbf v/d^4\) as a working assumption. The microscopic target
is to derive the low-frequency force-response tensor, including its
separation and finite-size dependence, from specified dynamics and
a specified stationary reference state.

\paragraph{Conditional Kubo formulation.}
Let \(\mathbf r\) denote the relative separation and
\(H(\mathbf r)\) a microscopic Hamiltonian. At fixed \(\mathbf r\),
choose a stationary reference state and define the connected
generalized-force operator
\[
\delta\hat F_i
=
-\frac{\partial H}{\partial r_i}
-\left\langle-\frac{\partial H}{\partial r_i}\right\rangle_{\mathbf r}.
\]
For a perturbation
\(\delta H(t)=-\delta r_j(t)\delta\hat F_j\), define
\[
\chi^R_{ij}(t;\mathbf r,\ell_A,\ell_B)
=
\frac{i}{\hbar}\theta(t)
\left\langle[\delta\hat F_i(t),\delta\hat F_j(0)]\right\rangle_{\mathbf r},
\qquad
\chi^R_{ij}(\omega)
=
\int_{-\infty}^{\infty} e^{i\omega t}
\chi^R_{ij}(t)\,\mathrm dt .
\]
For a reciprocal passive reference state, if the low-frequency
limit exists, define
\[
\Gamma_{ij}(\mathbf r;\ell_A,\ell_B)
=
\lim_{\omega\downarrow0}
\frac{\operatorname{Im}\chi^R_{(ij)}
(\omega;\mathbf r,\ell_A,\ell_B)}{\omega},
\qquad
\chi^R_{(ij)}=\frac{\chi^R_{ij}+\chi^R_{ji}}{2}.
\]
Passivity requires this dissipative tensor to be positive
semidefinite. The corresponding low-frequency force contribution is
\[
\delta F_i^{\mathrm{diss}}
=
-\Gamma_{ij}(\mathbf r;\ell_A,\ell_B)\dot r_j.
\]
This is a conditional response prescription, not a calculation of
the NAMCG force correlator, its reference state, or its passivity.

If the reference state and environment are rotationally invariant,
the response at nonzero separation can have distinct longitudinal
and transverse components,
\[
\Gamma_{ij}
=
\Gamma_L(d;\ell_A,\ell_B)\hat r_i\hat r_j
+\Gamma_T(d;\ell_A,\ell_B)
 \bigl(\delta_{ij}-\hat r_i\hat r_j\bigr).
\]
A radial response calculation determines only \(\Gamma_L\).
The isotropic scalar drag used in the laboratory requires
\(\Gamma_L=\Gamma_T\), or an explicit restriction to radial motion.

For example, the causal response family
\[
\chi^R_{ij,\tau}(\omega)
=
\frac{\chi_{0,ij}}{1-i\omega\tau},
\qquad \tau>0,
\]
with real symmetric positive-semidefinite \(\chi_{0,ij}\), has a
pole at \(\omega=-i/\tau\) and satisfies
\[
\chi^R_{ij,\tau}(0)=\chi_{0,ij},
\qquad
\lim_{\omega\downarrow0}
\frac{\operatorname{Im}\chi^R_{ij,\tau}(\omega)}{\omega}
=\tau\chi_{0,ij}.
\]
Thus a fixed static susceptibility does not determine the
low-frequency dissipative coefficient. This example is not a
NAMCG derivation.

Positivity of mutual information establishes \(I(A:B)\geq0\) for
each state. It does not establish an ordering between adiabatic
and non-adiabatic mutual informations, or the sign or separation
dependence of a dissipative response.

\paragraph{Static tail, distinct from the response.}
Chapters~\ref{ch:entanglement-equilibrium} and \ref{ch:dissipation}
and Appendix~\ref{app:rar-bridge} use the same large-separation
scaling,
\[
I(A:B)\sim c_{p_*}
\left(\frac{\ell_A\ell_B}{d^2}\right)^{p_*},
\qquad d\gg\ell_A,\ell_B,
\]
which is dimensionless. It is established in the literature for
R\'enyi mutual information of regions in a conformal vacuum, with the
von Neumann limit shown for free-scalar spheres
\cite{Cardy2013MI}; its leading power for NAMCG cores is an
assumption. Neither this static tail nor a coincidence of its power
with the assumed drag exponent determines the retarded force
response, its tensor structure, or its coefficient.

\paragraph{Missing step.}
Specify the microscopic dynamics, stationary reference state, and
force operators for the finite cores. Compute
\(\chi^R_{ij}(\omega;\mathbf r,\ell_A,\ell_B)\), justify the
passivity and low-frequency assumptions, and derive the
longitudinal and transverse responses. Determine whether they
reduce to the isotropic \(d^{-4}\) laboratory law and calculate its
coefficient before estimating \(\eta\) at tidal or LIGO-band scales.
Problem~I remains open.

\paragraph{Status.}
Coefficient, not load-bearing. A derivation of the drag from
microscopic dynamics would be desirable; its absence does not kill
the framework.

\section{Problem J: spherical ECG solutions}
\label{sec:status-J}

For the explicit invariant \(\mathcal P\) in~\eqref{eq:ECG}, consider
the reduced action per unit time and angular volume
\[
L[N,f]=N r^2\bigl(R-2\Lambda_{\mathrm{eff}}+\kappa\mathcal P\bigr),
\qquad
\kappa=16\pi G_{\mathrm{eff}}\beta_{\mathrm{ECG}}.
\]
This reduction assumes that the static, spherically symmetric
variations capture the corresponding field equations; an appropriate
symmetric-criticality argument is not established here. The reduced
Euler--Lagrange operators, including all higher-derivative terms, are
\[
\mathcal E_q
=\frac{\partial L}{\partial q}
-\frac{\mathrm d}{\mathrm dr}\frac{\partial L}{\partial q'}
+\frac{\mathrm d^2}{\mathrm dr^2}\frac{\partial L}{\partial q''}
-\cdots,\qquad q=N,f.
\]
Variations are taken with the fields and the derivatives required by
the order of \(L\) fixed at the endpoints (or with compact support);
otherwise the associated boundary terms must be retained. For the
specified \(\mathcal P\), these reduced equations have been checked
symbolically:
\[
\left.\mathcal E_f\right|_{N=1}\equiv0,\qquad
\left.\mathcal E_N\right|_{N=1}
=\frac{\mathrm d}{\mathrm dr}
\left[2r(1-f)+3\kappa\mathcal B[f]
-\frac23\Lambda_{\mathrm{eff}}r^3\right].
\]
Thus \(N=1\) is a consistent vacuum branch of the reduced equations,
and, with \(\alpha=-\kappa\), the lapse equation has first integral
\[
\mathcal F[f]-\frac23\Lambda_{\mathrm{eff}}r^3=4M,\qquad
\mathcal F[f]=2r(1-f)-3\alpha\,\mathcal B[f],
\]
where \(\alpha\) is the reduced ECG coupling and
\[
\mathcal B[f]
=
\frac43(f')^3
-\frac8r\,ff''
 \left(\frac{rf'}2+1-f\right)
+\frac4r(f')^2
-\frac8{r^2}ff'(f-1).
\]
For the rational profile with \(a=r_0^3\) and
\(\Lambda_{\mathrm{eff}}=0\),
\[
f_{\mathrm{rat}}(r)
=
1-\frac{2Mr^2}{r^3+a}
=
1-\frac{2M}{r}+\frac{2Ma}{r^4}+O(r^{-7}),
\]
substitution gives
\[
\mathcal F[f_{\mathrm{rat}}]-4M
=-\frac{4Ma}{r^3}-\frac{432\alpha M^2}{r^5}+O(r^{-6}).
\]
The \(r^{-3}\) term cannot be cancelled by the cubic contribution,
which begins at \(r^{-5}\). Hence, for the displayed equation, this
profile is not an exact vacuum solution when \(Mr_0\ne0\).

Moreover, a smooth regular center has \(f(r)=1+O(r^2)\), which
implies \(\mathcal B[f]\to0\) and
\(\mathcal F(0)-\tfrac23\Lambda_{\mathrm{eff}}0^3=0\).
Consequently the integration constant vanishes: this displayed
equation has no smooth regular-centre vacuum solution with a nonzero
integration constant, for any \(\Lambda_{\mathrm{eff}}\). For
\(\Lambda_{\mathrm{eff}}=0\), the asymptotic falloff
\(f=1-2M/r+o(r^{-1})\) identifies that constant as \(4M\) and \(M\)
as the ADM mass. For nonzero \(\Lambda_{\mathrm{eff}}\), its
interpretation is not automatically the ADM mass.

\paragraph{Maximally symmetric branch.}
For \(f=1-kr^2\), the first integral gives
\[
\frac{\mathcal F[f]}{r^3}=2k-16\alpha k^3,\qquad
k-8\alpha k^3=\frac{\Lambda_{\mathrm{eff}}}{3}.
\]
This is a conditional background-curvature relation, not a stability
claim.

\paragraph{Large-radius tail.}
For the displayed equation, write \(f=1-2M/r+\delta(r)\). Then
\(2r(1-f)=4M-2r\delta\) and
\(\mathcal B[1-2M/r]=144M^2/r^5+O(r^{-6})\). The leading
correction and the next asymptotic order give
\[
f(r)=1-\frac{2M}{r}-\frac{216\,\alpha M^2}{r^6}
+\frac{368\,\alpha M^3}{r^7}+\cdots .
\]
The rational profile instead has a \(+2Ma/r^4\) term. The tail
expansion is an asymptotic check of the displayed first integral; it
does not determine the model parameter \(r_0\) or establish a match to
an independently published expansion. In the stated action
normalization \(\alpha=-16\pi G_{\mathrm{eff}}\beta_{\mathrm{ECG}}\).
If \(G_{\mathrm{eff}}>0\), then \(\alpha>0\) iff
\(\beta_{\mathrm{ECG}}<0\), which fixes the sign of the leading
correction. For \(M>0\), its magnitude defines the characteristic
scale \(r_*\sim|\alpha M|^{1/5}\); it is not a core scale.

\paragraph{Status.}
Conditional on the stated symmetric reduction, the reduced-action
identities, first integral, regular-centre constraint, maximally
symmetric curvature relation, and asymptotic tail above have been
checked for the specified \(\mathcal P\). The cited Einstein-like
linearized spectrum on maximally symmetric backgrounds is a separate
result; no nonlinear, sourced, or non-maximally-symmetric stability
claim follows. Problem~J
remains open because the reduction's relation to the full field
equations and the cubic truncation and its matching to the matrix
theory are not established. A regular massive core still requires a
specified source, different boundary conditions, or modified
field equations; the entanglement stress for such a source is not specified.

\section{What this appendix does not do}
\label{sec:status-not}

It does not produce an unsuppressed laboratory coupling between a
controllable external field and \(T_{\mu\nu}^{\mathrm{ent}}\) or
\(G_{\mathrm{eff}}\). The observational window therefore remains empty.
It does not alter the chain of Chapter~\ref{ch:falsifiers} or the agenda of
Chapter~\ref{ch:agenda}.